\subsection{Sufijos por Responsabilidad}

\subsubsection{Regla: sufijos definidos por .NET}

\textbf{Regla:} Se aplican los sufijos establecidos por las guías de diseño de .NET: \texttt{Attribute} para atributos, \texttt{EventArgs} para datos de eventos, \texttt{EventHandler} para delegados de eventos, \texttt{Callback} para otros delegados, \texttt{Exception} para excepciones, \texttt{Dictionary} para diccionarios y \texttt{Collection} para colecciones personalizadas (Microsoft, 2023c).

\textbf{Código correcto:}
\begin{lstlisting}[style=csharpstyle]
public sealed class TurnChangedEventArgs : EventArgs
{
    public TurnChangedEventArgs(Player currentPlayer)
    {
        CurrentPlayer = currentPlayer;
    }

    public Player CurrentPlayer { get; }
}
\end{lstlisting}

\textbf{Código incorrecto:}
\begin{lstlisting}[style=csharpstyle]
public sealed class TurnChangedData : EventArgs
{
    public TurnChangedData(Player currentPlayer)
    {
        CurrentPlayer = currentPlayer;
    }

    public Player CurrentPlayer { get; }
}
\end{lstlisting}

\subsubsection{Regla: sufijos arquitectónicos del proyecto}

\textbf{Regla:} Para expresar una responsabilidad arquitectónica real, el proyecto adopta los sufijos \texttt{Service}, \texttt{Repository}, \texttt{Controller}, \texttt{Factory}, \texttt{Provider} y \texttt{Validator}. Se prohíbe añadirlos cuando la clase no desempeña esa responsabilidad.

\textbf{Código correcto:}
\begin{lstlisting}[style=csharpstyle]
public sealed class SaveGameService
{
    public void Save(GameState state)
    {
        Persist(state);
    }
}
\end{lstlisting}

\textbf{Código incorrecto:}
\begin{lstlisting}[style=csharpstyle]
public sealed class SaveGameManager
{
    public void Save(GameState state)
    {
        Persist(state);
    }
}
\end{lstlisting}
